% !PRINTED 482
\weberpage{482}

A system of three functions which undergo ternary substitution can be formed in the following way\footnote{F.\ \gesperrt{Klein}, \glqq Ueber die Auflösung gewisser Gleichungen vom 7\textsuperscript{ten} und 8\textsuperscript{ten} Grade\grqq. Mathem. Annalen, Bd. XV (1879).}.

We first introduce a system of auxiliary variables $x_1, x_2, x_3$, which we subject to the substitutions of the group $G_{168}$, and from these we form the function $z_0$ belonging to the group $G_{24}$, together with its conjugates $z_r$ [\S.\,125, (7)]:
\weq{2}{z_r = \varepsilon^{2r}x_1^2 + \varepsilon^{4r}x_2^2 + \varepsilon^r x_3^2 - \frac{1-\sqrt{-7}}{2}(\varepsilon^{-r}x_2x_3 + \varepsilon^{-2r}x_3x_1 + \varepsilon^{-4r}x_1x_2).}

To this we now take a function $\eta_0$ of the roots $\lambda_r$ of our equation of the 7th degree which belongs to the group $P_{24}$, e.\,g.\ a rational function of $\lambda_0$, and the conjugate values $\eta_0, \eta_1, \ldots, \eta_6$, and form the sum
\weq{3}{\psi = \eta_0 z_0 + \eta_1 z_1 + \eta_2 z_2 + \eta_3 z_3 + \eta_4 z_4 + \eta_5 z_5 + \eta_6 z_6,}
which is a quadratic function of the $x$, whose coefficients depend on the roots $\lambda_r$.

\gesperrt{This function $\psi$ does not change when the variables $(x)$ are subjected to a substitution of the group $G_{168}$ and the roots $\lambda_r$ at the same time to the corresponding permutation from $P_{168}$.} For by this simultaneous operation only the summands in the sum (2) are interchanged among one another, and hence the sum itself is not changed.

We may therefore designate $\psi$ as a \gesperrt{simultaneous invariant} of the groups $G_{168}$ and $P_{168}$.

If we arrange the function $\psi$ according to the variables $x_1, x_2, x_3$, we obtain
\weq{4}{\psi = p_1x_1^2 + p_2x_2^2 + p_3x_3^2 + 2q_1x_2x_3 + 2q_2x_3x_1 + 2q_3x_1x_2,}
where for abbreviation
\weq{5}{\begin{aligned}
p_1 &= \sum^r \varepsilon^{2r}\eta_r, & q_1 &= -\frac{1-\sqrt{-7}}{4}\sum^r \varepsilon^{-r}\eta_r\\
p_2 &= \sum^r \varepsilon^{4r}\eta_r, & q_2 &= -\frac{1-\sqrt{-7}}{4}\sum^r \varepsilon^{-2r}\eta_r\\
p_3 &= \sum^r \varepsilon^{r}\eta_r, & q_3 &= -\frac{1-\sqrt{-7}}{4}\sum^r \varepsilon^{-4r}\eta_r
\end{aligned}}
